% ======================================================================
\chapter[Фундаментальные частицы и взаимодействия. Стандартная модель]%
        {Фундаментальные частицы и взаимодействия. Стандартная модель}
\label{t:14}
% ======================================================================

\begin{qbox}
\textbf{Билет 14.} Фундаментальные частицы и фундаментальные взаимодействия (Стандартная модель:
лептоны, кварки и переносчики взаимодействий, бозон Хиггса).
\end{qbox}

\section{Содержание билета}

\textbf{Подтемы:}
\begin{enumerate}
 \item понятие фундаментальной частицы; фермионы вещества и бозоны-переносчики;
 \item четыре фундаментальных взаимодействия: интенсивность, радиус действия, переносчики;
 \item обменный механизм взаимодействия и оценка радиуса (Юкава); диаграммы Фейнмана;
 \item лептоны (три поколения), кварки (ароматы, цвет, дробные заряды);
 \item переносчики: фотон, $W^\pm$, $Z^0$, глюоны; бозон Хиггса и механизм генерации масс;
 \item законы сохранения; границы применимости Стандартной модели.
\end{enumerate}

\textbf{Необходимые понятия:} спин, бозоны и фермионы (\tref{t:12}, \tref{t:13}); соотношение
неопределённостей энергия--время (\tref{t:02}); релятивистская связь $E^2=p^2c^2+m^2c^4$; античастицы.

\section{Определения и физические принципы}

\begin{definition}[фундаментальная частица]
Частица, у которой на достигнутых масштабах (до $\sim10^{-19}$~м) не обнаружено внутренней структуры.
В Стандартной модели (СМ) фундаментальны 12 фермионов (6 лептонов и 6 кварков) с их античастицами,
12 калибровочных бозонов (фотон, $W^+$, $W^-$, $Z^0$, 8 глюонов) и бозон Хиггса.
\end{definition}

\begin{definition}[фундаментальное взаимодействие]
Один из четырёх известных типов взаимодействий, к которым сводятся все силы в природе: сильное,
электромагнитное, слабое, гравитационное. В квантовой теории поля взаимодействие осуществляется
обменом квантами поля — частицами-переносчиками (бозонами со спином 1; гравитация — гипотетический
гравитон со спином 2).
\end{definition}

\begin{definition}[античастица]
Частица с той же массой, спином и временем жизни, но противоположными зарядами (электрическим,
барионным, лептонным, цветом). Некоторые нейтральные бозоны ($\gamma$, $Z^0$, $H$, $\pi^0$) совпадают со своими
античастицами.
\end{definition}

\emph{Область применимости СМ.} СМ — квантовая теория поля с калибровочной группой
$SU(3)_C\times SU(2)_L\times U(1)_Y$, описывающая сильное, слабое и электромагнитное взаимодействия.
Гравитация в неё не входит. СМ не объясняет массы нейтрино (в исходной версии нейтрино безмассовы),
тёмную материю, барионную асимметрию Вселенной; значения масс и констант связи — свободные параметры.

\section{Основная теория}

\subsection{Проблема: от «зоопарка» частиц к упорядоченной схеме}

К 1960-м годам было открыто более сотни «элементарных» частиц — адронов. Их классификация
(Гелл-Ман, Цвейг, 1964) показала, что адроны составные — построены из кварков (\tref{t:15}).
Одновременно развивалась калибровочная теория взаимодействий: квантовая электродинамика (КЭД),
объединённая электрослабая теория (Глэшоу, Вайнберг, Салам, 1967--68) и квантовая хромодинамика
(КХД, 1973). Вместе они составили Стандартную модель, все предсказанные частицы которой открыты
(последней — бозон Хиггса, 2012).

\subsection{Фундаментальные взаимодействия}

\begin{center}
\small
\begin{tabularx}{\linewidth}{@{}lccc>{\raggedright\arraybackslash}X@{}}
\toprule
Взаимодействие & Переносчик & Масса переносчика & Радиус & Отн. интенсивность (при $\sim1$~ГэВ) \\
\midrule
Сильное & 8 глюонов $g$ & 0 & $\sim1$~фм (конфайнмент) & $\alpha_s\sim1$ \\
Электромагнитное & фотон $\gamma$ & 0 & $\infty$ & $\alpha=1/137$ \\
Слабое & $W^\pm$, $Z^0$ & 80,4 и 91,2 ГэВ & $\sim2\cdot10^{-3}$~фм & $G_Fm_p^2/(\hbar c)^3\approx10^{-5}$ \\
Гравитационное & (гравитон) & 0 & $\infty$ & $Gm_p^2/\hbar c\approx6\cdot10^{-39}$ \\
\bottomrule
\end{tabularx}
\end{center}

\noindent Безразмерные «константы связи» характеризуют силу взаимодействия: $\alpha=e^2/\hbar c$ для
электромагнитного, $\alpha_s=g_s^2/\hbar c$ для сильного. Для слабого при низких энергиях роль константы
играет $G_F$ (размерная), и интенсивность растёт с энергией как $G_FE^2$; при $E\gtrsim M_Wc^2$ слабое и
электромагнитное взаимодействия сравнимы по силе. Характерные времена распадов: сильные — $10^{-23}$~с,
электромагнитные — $10^{-16}$--$10^{-20}$~с, слабые — $\gtrsim10^{-13}$~с (вплоть до 880~с для нейтрона).

\paragraph{Какие частицы в чём участвуют.} В сильном — только кварки и глюоны (а значит, адроны);
в электромагнитном — заряженные частицы; в слабом — все фермионы (включая нейтрино, которые участвуют
\emph{только} в слабом и гравитационном); в гравитационном — все.

\subsection{Обменный механизм. Радиус взаимодействия}

Взаимодействие двух частиц описывается как испускание виртуального кванта одной частицей и
поглощение другой. Испускание частицы массы $M$ покоящейся частицей нарушает закон сохранения энергии
на $\Delta E\gtrsim Mc^2$; по соотношению неопределённостей такое состояние может существовать время
$\Delta t\sim\hbar/Mc^2$, за которое квант пролетает расстояние (вывод~\ref{der:14:yukawa})
\begin{equation}\label{eq:14:range}
  r_0\sim c\Delta t=\frac{\hbar}{Mc}.
\end{equation}
Это комптоновская длина переносчика: безмассовый переносчик даёт дальнодействие (закон $1/r^2$), массивный —
короткодействие с потенциалом Юкавы
\begin{equation}\label{eq:14:yukawa}
  U(r)\propto-\frac{g^2}{r}\ee^{-r/r_0}.
\end{equation}
Х.~Юкава (1935) по радиусу ядерных сил $\sim1{,}4$~фм предсказал переносчик с массой $\sim140$~МэВ — пи-мезон
(открыт в 1947). Для $W$-бозона $r_0=\hbar c/M_Wc^2=197\MeV\cdot\text{фм}/80\,400\MeV\approx2{,}5\cdot10^{-3}$~фм.

\paragraph{Диаграммы Фейнмана} — графическая запись амплитуд процессов: сплошные линии — фермионы,
волнистые — фотоны и $W$, $Z$, спиральные — глюоны; вершина соответствует элементарному акту
взаимодействия, в котором сохраняются заряды; каждая вершина вносит в амплитуду множитель
$\sim\sqrt{\alpha}$ (константа связи). Примеры — на рис.~\ref{fig:14:feyn}.

\begin{figure}[htbp]
\centering
\begin{tikzpicture}[scale=0.85, photon/.style={decorate, decoration={snake, amplitude=1.1mm, segment length=2.6mm}, thick},
  gluon/.style={decorate, decoration={coil, amplitude=1.3mm, segment length=1.6mm, aspect=0.8}, thick},
  fermion/.style={thick, postaction={decorate}, decoration={markings, mark=at position 0.55 with {\arrow{Stealth[length=2mm]}}}}]
  % (a) e e рассеяние
  \begin{scope}
   \draw[fermion] (0,0) -- (1,1); \draw[fermion] (1,1) -- (0,2);
   \draw[fermion] (3,0) -- (2,1); \draw[fermion] (2,1) -- (3,2);
   \draw[photon] (1,1) -- (2,1);
   \node[left, font=\small] at (0,0) {$e^-$}; \node[left, font=\small] at (0,2) {$e^-$};
   \node[right, font=\small] at (3,0) {$e^-$}; \node[right, font=\small] at (3,2) {$e^-$};
   \node[above, font=\small] at (1.5,1.15) {$\gamma$};
   \node[font=\small] at (1.5,-0.7) {(а) ЭМ};
  \end{scope}
  % (b) q q рассеяние глюоном
  \begin{scope}[xshift=4.7cm]
   \draw[fermion] (0,0) -- (1,1); \draw[fermion] (1,1) -- (0,2);
   \draw[fermion] (3,0) -- (2,1); \draw[fermion] (2,1) -- (3,2);
   \draw[gluon] (1,1) -- (2,1);
   \node[left, font=\small] at (0,0) {$q$}; \node[left, font=\small] at (0,2) {$q$};
   \node[right, font=\small] at (3,0) {$q$}; \node[right, font=\small] at (3,2) {$q$};
   \node[above, font=\small] at (1.5,1.2) {$g$};
   \node[font=\small] at (1.5,-0.7) {(б) сильное};
  \end{scope}
  % (c) бета-распад нейтрона
  \begin{scope}[xshift=9.4cm]
   \draw[fermion] (0,0) -- (1,1); \draw[fermion] (1,1) -- (0,2);
   \draw[photon] (1,1) -- (2.2,1.3);
   \draw[fermion] (2.2,1.3) -- (2.6,2.3);
   \draw[fermion] (3.3,1.9) -- (2.2,1.3);
   \node[left, font=\small] at (0,0) {$d$}; \node[left, font=\small] at (0,2) {$u$};
   \node[above, font=\small] at (2.6,2.3) {$e^-$}; \node[right, font=\small] at (3.3,1.9) {$\bar\nu_e$};
   \node[below, font=\small] at (1.75,1.0) {$W^-$};
   \node[font=\small] at (1.5,-0.7) {(в) слабое};
  \end{scope}
\end{tikzpicture}
\caption{Диаграммы Фейнмана (время — снизу вверх): (а) рассеяние электронов обменом фотоном;
(б) рассеяние кварков обменом глюоном; (в) слабый распад $d\to u+e^-+\bar\nu_e$, лежащий в основе
бета-распада нейтрона (стрелка, направленная против времени, — античастица).}
\label{fig:14:feyn}
\end{figure}

\subsection{Фермионы: лептоны и кварки}

Фермионы со спином $\frac12$ образуют три поколения (семейства), одинаковые по зарядам и взаимодействиям,
но различающиеся массами.

\begin{table}[htbp]
\centering
\small
\caption{Фундаментальные фермионы Стандартной модели (массы — в единицах $c^{-2}$, приближённо).}
\label{tab:14:fermions}
\begin{tabular}{@{}lcccc@{}}
\toprule
& I поколение & II поколение & III поколение & Заряд $Q/e$ \\
\midrule
\multirow{2}{*}{Кварки}
 & $u$ — 2,2 МэВ & $c$ — 1,27 ГэВ & $t$ — 173 ГэВ & $+\tfrac23$ \\
 & $d$ — 4,7 МэВ & $s$ — 95 МэВ & $b$ — 4,18 ГэВ & $-\tfrac13$ \\
\midrule
\multirow{2}{*}{Лептоны}
 & $e^-$ — 0,511 МэВ & $\mu^-$ — 105,7 МэВ & $\tau^-$ — 1777 МэВ & $-1$ \\
 & $\nu_e$ & $\nu_\mu$ & $\nu_\tau$ & $0$ \\
\bottomrule
\end{tabular}
\end{table}

\paragraph{Лептоны} не участвуют в сильном взаимодействии. Заряженные лептоны $e,\mu,\tau$ участвуют в
электромагнитном и слабом взаимодействиях; $\mu$ и $\tau$ нестабильны ($\tau_\mu=2{,}2$~мкс, $\tau_\tau=2{,}9\cdot10^{-13}$~с).
Нейтрино — нейтральные, участвуют только в слабом взаимодействии; их массы ненулевые (доказано
наблюдением нейтринных осцилляций; Нобелевская премия 2015), но очень малы — меньше $\sim1$~эВ. Каждому
поколению приписывают лептонное число ($L_e$, $L_\mu$, $L_\tau$): $+1$ для лептона, $-1$ для антилептона.

\paragraph{Кварки} имеют дробный электрический заряд ($+\frac23e$ или $-\frac13e$), барионное число $\frac13$ и
дополнительное квантовое число — \emph{цвет} (три значения: «красный», «зелёный», «синий»). В свободном
состоянии не наблюдаются (конфайнмент): существуют только бесцветные комбинации — адроны (\tref{t:15}).
Шесть типов кварков называют \emph{ароматами}: $u,d,s,c,b,t$. Массы кварков указаны условно: лёгкие кварки
нельзя «взвесить» в свободном состоянии; их массы определяются косвенно.

\subsection{Переносчики взаимодействий}

\begin{itemize}
 \item \textbf{Фотон} $\gamma$: спин 1, масса 0, заряда нет. Взаимодействует с электрическим зарядом.
 \item \textbf{$W^\pm$, $Z^0$}: спин 1, массы $80{,}4$ и $91{,}2$~ГэВ; открыты в ЦЕРН в 1983~г. (К.~Руббиа, С.~ван~дер~Мер).
       $W^\pm$ меняют аромат фермиона (заряженные токи), $Z^0$ — нет (нейтральные токи, открыты в 1973~г.).
 \item \textbf{Глюоны} $g$: спин 1, масса 0, электрически нейтральны, но несут цвет (8 цветовых комбинаций);
       поэтому глюоны взаимодействуют друг с другом — это принципиальное отличие КХД от КЭД (\tref{t:15}).
\end{itemize}

\subsection{Бозон Хиггса}

Калибровочная симметрия требует, чтобы переносчики были безмассовыми, а слабые бозоны массивны.
Противоречие разрешается \emph{механизмом Хиггса} (Браут, Энглер, Хиггс, 1964): пространство заполнено
скалярным полем, имеющим ненулевое вакуумное среднее $v\approx246$~ГэВ. Взаимодействие с ним даёт массы:
\begin{itemize}
 \item $W^\pm$ и $Z^0$ — через спонтанное нарушение электрослабой симметрии ($M_W=\frac12gv$, $M_Z=M_W/\cos\theta_W$);
       фотон остаётся безмассовым;
 \item фермионам — через юкавские взаимодействия: $m_f=y_fv/\sqrt2$, поэтому константа связи фермиона с
       бозоном Хиггса пропорциональна его массе.
\end{itemize}
Квант этого поля — \emph{бозон Хиггса} $H$: спин 0, электрически нейтрален, масса $\approx125$~ГэВ. Открыт в
2012~г. на Большом адронном коллайдере (эксперименты ATLAS и CMS); Нобелевская премия 2013~г.
(Энглер, Хиггс). Он распадается преимущественно в пары $b\bar b$, а также $WW^*$, $\gamma\gamma$, $ZZ^*$ и др.
Механизм Хиггса не объясняет основную часть массы обычного вещества: масса протона ($938$~МэВ)
определяется энергией глюонного поля и движения кварков (КХД), а не массами $u$ и $d$ кварков ($\sim10$~МэВ в сумме).

\subsection{Законы сохранения}

\begin{center}
\small
\begin{tabular}{@{}lcccc@{}}
\toprule
Величина & Сильное & ЭМ & Слабое \\
\midrule
энергия, импульс, момент импульса, электрический заряд & да & да & да \\
барионное число $B$ & да & да & да \\
лептонные числа $L_e,L_\mu,L_\tau$ & да & да & да$^{*}$ \\
аромат (странность $S$, очарование, \dots) & да & да & нет \\
пространственная чётность $P$ & да & да & нет \\
\bottomrule
\end{tabular}

{\footnotesize $^{*}$ В процессах взаимодействия; нейтринные осцилляции нарушают сохранение отдельных $L_e,L_\mu,L_\tau$,
сохраняя полное $L$.}
\end{center}
Пример применения: распад $p\to e^++\gamma$ разрешён зарядом и энергией, но запрещён сохранением $B$ и $L$;
экспериментально $\tau_p>10^{34}$~лет.

\section{Математические выводы}

\begin{derivation}[label=der:14:yukawa]{радиус обменного взаимодействия}
\emph{Допущения:} оценка по порядку величины; покоящийся нуклон (или другая частица) испускает
виртуальный квант массы $M$.

В начальном состоянии энергия $E_0=m_Nc^2$. В промежуточном состоянии есть нуклон и квант с энергией
не менее $Mc^2$, т.\,е. закон сохранения энергии нарушен на $\Delta E\ge Mc^2$. Соотношение неопределённостей
$\Delta E\,\Delta t\sim\hbar$ разрешает такое состояние на время $\Delta t\sim\hbar/Mc^2$. Даже двигаясь со скоростью,
близкой к $c$, квант удаляется не более чем на
\[
  r_0\sim c\Delta t=\frac{\hbar}{Mc}.
\]
\emph{Размерность:} $[\hbar/Mc]=\text{эрг}\cdot\text{с}/(\text{г}\cdot\text{см}/\text{с})=\text{см}$.
\emph{Численно:} $\hbar c=197$~МэВ$\cdot$фм; для $\pi$-мезона ($139{,}6$~МэВ) $r_0=1{,}41$~фм — радиус ядерных
сил; для $W$ ($80\,400$~МэВ) — $2{,}5\cdot10^{-3}$~фм; для фотона ($M=0$) $r_0=\infty$.

\emph{Потенциал Юкавы.} Статическое поле массивного скалярного кванта удовлетворяет уравнению
$(\nabla^2-1/r_0^2)\phi=-4\pi g\delta(\vect r)$ (аналог уравнения Пуассона), решение которого
$\phi=g\,\ee^{-r/r_0}/r$. При $r_0\to\infty$ получаем кулоновский потенциал.
\end{derivation}

\begin{derivation}[label=der:14:compare]{сравнение интенсивностей взаимодействий}
\emph{Электромагнитное против гравитационного} для двух протонов:
\[
  \frac{F_{\text{эм}}}{F_{\text{гр}}}=\frac{e^2}{Gm_p^2}=\frac{(4{,}8\cdot10^{-10})^2}{6{,}67\cdot10^{-8}\cdot(1{,}67\cdot10^{-24})^2}\approx1{,}2\cdot10^{36}.
\]
Отношение не зависит от расстояния (обе силы $\propto1/r^2$). Гравитация в микромире пренебрежимо мала;
в макромире она доминирует, потому что тела электрически нейтральны, а «гравитационный заряд» (масса) всегда
положителен и складывается.

\emph{Слабое против электромагнитного.} Безразмерная мера слабого взаимодействия при энергии $E$:
$G_FE^2/(\hbar c)^3$. При $E\sim m_pc^2\approx1$~ГэВ это $\approx1{,}2\cdot10^{-5}$, тогда как $\alpha\approx7\cdot10^{-3}$. Константа
слабой связи $g$ сама по себе не мала ($g^2/4\pi\hbar c\approx1/30$); малость обусловлена множителем
$(E/M_Wc^2)^2$ — слабость «кажущаяся», связанная с массивностью переносчика:
$G_F/\sqrt2=g^2/8M_W^2$ (в системе $\hbar=c=1$).
\end{derivation}

\begin{derivation}[label=der:14:charges]{проверка зарядов в Стандартной модели}
Сумма электрических зарядов всех фермионов одного поколения с учётом трёх цветов кварков:
\[
  \underbrace{(-1)}_{e}+\underbrace{0}_{\nu}+3\Bigl(\underbrace{\tfrac23}_{u}+\underbrace{(-\tfrac13)}_{d}\Bigr)=-1+3\cdot\tfrac13=0 .
\]
Это условие сокращения аномалий, необходимое для внутренней согласованности (перенормируемости) теории;
оно связывает число цветов с дробностью зарядов кварков и указывает на глубокую связь лептонов и кварков.
\end{derivation}

\section{Итоговые формулы}

\keyformula{r_0\sim\frac{\hbar}{Mc},\qquad U(r)\propto-\frac{g^2}{r}\ee^{-r/r_0}}
 {$M$ — масса переносчика, $g$ — константа связи, $r_0$ — радиус действия.}
 {оценка по порядку величины; статический предел.}
 {радиус взаимодействия по массе переносчика и наоборот (пион, $W$).}

\keyformula{\alpha=\frac{e^2}{\hbar c}\approx\frac{1}{137},\quad \alpha_s\sim1\ (\text{при}\sim1\ \text{ГэВ}),\quad \frac{G_F}{(\hbar c)^3}=1{,}166\cdot10^{-5}\ \text{ГэВ}^{-2}}
 {$\alpha$, $\alpha_s$ — константы электромагнитной и сильной связи; $G_F$ — постоянная Ферми.}
 {значения зависят от энергии (бегущие константы).}
 {сравнение интенсивностей взаимодействий и времён процессов.}

\keyformula{m_f=\frac{y_fv}{\sqrt2},\quad M_W=\frac{gv}{2},\quad M_Z=\frac{M_W}{\cos\theta_W},\quad v\approx246\ \text{ГэВ}}
 {$y_f$ — юкавская константа фермиона, $g$ — константа $SU(2)$, $\theta_W$ — угол Вайнберга ($\sin^2\theta_W\approx0{,}23$).}
 {Стандартная модель, $\hbar=c=1$ (для справки).}
 {происхождение масс $W$, $Z$ и фермионов.}

\section{Физическая интерпретация}

\textbf{Структура.} Вещество состоит из фермионов первого поколения ($u$, $d$, $e$); второе и третье поколения
нестабильны и рождаются в столкновениях высоких энергий. Число поколений с лёгкими нейтрино равно трём — это
установлено по ширине распада $Z^0$ (LEP, 1989--2000): $N_\nu=2{,}98\pm0{,}01$.

\textbf{Масса и радиус.} Связь $r_0=\hbar/Mc$ объясняет, почему дальнодействующие силы (электромагнитная,
гравитационная) имеют безмассовых переносчиков. Сильное взаимодействие короткодействующее по другой причине:
глюоны безмассовы, но из-за конфайнмента цветные объекты не распространяются дальше $\sim1$~фм.

\textbf{Экспериментальный статус.} Предсказания СМ подтверждены с высокой точностью (аномальный магнитный момент
электрона, свойства $W$ и $Z$, сечения процессов КХД). За её пределами остаются гравитация, тёмная материя и
тёмная энергия, массы нейтрино и преобладание вещества над антивеществом.

\section{Что необходимо знать на экзамене}

\begin{examcheck}
\checkhead{Определения}
\begin{itemize}
 \item фундаментальная частица; лептоны, кварки (ароматы, цвет), поколения; переносчики; античастицы;
 \item четыре взаимодействия: кто участвует, переносчики, радиусы, относительная сила.
\end{itemize}
\checkhead{Формулы и числа}
\begin{itemize}
 \item $r_0\sim\hbar/Mc$; потенциал Юкавы; $\hbar c=197$~МэВ$\cdot$фм;
 \item $\alpha\approx1/137$, $\alpha_s\sim1$, $G_F$; $M_W\approx80$~ГэВ, $M_Z\approx91$~ГэВ, $M_H\approx125$~ГэВ;
 \item заряды кварков $+\frac23$, $-\frac13$; массы $e,\mu,\tau$.
\end{itemize}
\checkhead{Выводы}
\begin{itemize}
 \item оценка радиуса обменного взаимодействия (пион, $W$-бозон);
 \item сравнение гравитационного и электромагнитного взаимодействий; «кажущаяся» слабость слабого взаимодействия.
\end{itemize}
\checkhead{Интерпретации}
\begin{itemize}
 \item роль бозона Хиггса (массы $W$, $Z$, фермионов) и почему масса протона — не от Хиггса;
 \item законы сохранения ($B$, $L$, аромат, чётность) и их нарушение слабым взаимодействием; ограничения СМ.
\end{itemize}
\end{examcheck}
